\section{The obstruction to a global finite construction}
\label{obs:section}

The object $U$ of Proposition~\ref{glue:object} has finite homotopy
groups in every degree, and its completion is the finite construction
\[
 Y=\bigoplus_{r=0}^{D-1}m[r],\qquad
 m=\cofib\bigl(a:\widehat T'\longrightarrow\widehat T'\bigr).
\]
The remaining condition in the transfer criterion is
$U\notin\Thick(T')$.  We must rule out every finite tower of shifts
of $T'$ and every retract of such a tower.  Ruling out a retract
of a finite sum of shifts would still leave longer towers untreated.

We first isolate the finite-tower argument.  Its input is an exact
functor to perfect modules over a periodic algebra over a domain
with principal left and right ideals.  The criterion uses ordinary
homology and the homology of module duals
to pass backward through any finite tower: the dual test removes
the free part of a tower tail, and the ordinary test then uses
torsion to force vanishing.  We next construct a flat extension of
the operation coefficients for which $Y$ has nonzero torsion
homology.  The final part verifies the relation between the two
tests for maps that come from global objects.

\subsection{A criterion for arbitrary finite constructions}

We state the two tests abstractly before constructing the
coefficient extension that will make them apply to the operator
$a$ and the gluing lattice.  Torsion homology by itself
cannot be the obstruction: the cone of $2:\Z\to\Z$ is a two-cell
object with homology $\Z/2$.  The additional information will
restrict the maps that exist in the original category.

Let $\mathscr B_0=S_0[t^{\pm1}]$, where $S_0$ is a domain with principal
left and right ideals, $t$ is central of positive even homological
degree $D_0$, and the differential is zero.  Write $\Perf(\mathscr B_0)$
for the full subcategory of the derived category consisting of
perfect right dg $\mathscr B_0$-modules.
For a right module over a domain, we call an element \emph{torsion}
if it is annihilated on the right by a nonzero ring element.  In the
principal domains used below the Ore property makes these elements
a submodule.  A torsion module has no nonzero map to a free module:
the image of a torsion element would be annihilated by a nonzero
element of the domain.

For $M\in\Perf(\mathscr B_0)$ its module dual is the left dg module
\[
 M^*=R\Hom_{\mathscr B_0}(M,\mathscr B_0).
\]
A map $f:M\to N$ induces
$H(f):H(M)\to H(N)$ and
$H(f^*):H(N^*)\to H(M^*)$.  The dual test has the opposite
direction, and in both tests vanishing means zero in every degree.

Testing finite constructions by maps invisible to homology is part
of the theory of ghosts and projective classes developed by
Christensen~\cite[Theorem~3.5 and Section~8]{Christensen1998}.
The following criterion adds a relation between the two tests for
maps in the original category.  Its proof is included.

\begin{theorem}[A finite-tower criterion]
\label{rew:obs:tower-criterion}
Let $\mathcal C$ be an idempotent-complete triangulated category,
let $T_0,U_0\in\mathcal C$, and let
\[
 \Psi:\mathcal C\longrightarrow\Perf(\mathscr B_0)
\]
be an exact functor with $\Psi(T_0)\simeq\mathscr B_0$.  Assume that
$H(\Psi(U_0))$ is torsion as a right $S_0$-module in every degree
and is nonzero in at least one degree.  Assume also that, for every
$P\in\Thick_{\mathcal C}(T_0)$ and every map $g:U_0\to P$ in
$\mathcal C$,
\begin{equation}
\label{rew:obs:two-tests}
 H(\Psi(g))=0
 \quad\Longrightarrow\quad
 H\bigl((\Psi(g))^*\bigr)=0.
\end{equation}
Then $U_0\notin\Thick_{\mathcal C}(T_0)$.
\end{theorem}

Here $\Thick_{\mathcal C}(T_0)$ is the ordinary thick closure,
using isomorphisms, finite sums, all integer shifts, triangles
and retracts.
The quantifier in \eqref{rew:obs:two-tests} concerns maps in
$\mathcal C$.  In our application, their global origin is precisely
what imposes a constraint on the gluing lattice.

The proof uses two algebraic facts.  The first separates free
homology from torsion homology over the coefficient ring.

\begin{lemma}[Periodic splitting]
\label{obs:periodic-splitting}
Let $S_0$ be a domain whose left and right ideals are principal, and
let $D_0$ be a positive even integer.  Every finitely generated right
$S_0$-module is a direct sum of a torsion module and a finite free
module and has a finite free resolution of length at most one.
Every perfect right dg module over
$S_0[t^{\pm1}]$, with $|t|=D_0$ and zero differential, splits in its
derived category into the $D_0$-periodizations of its shifted homology
modules.  In particular it splits as a finite sum of free shifts and
a perfect module with torsion homology.
\end{lemma}

\begin{proof}
First, $S_0$ is noetherian on both sides, since every one-sided
ideal is principal.  Nonzero right ideals $dS_0$ and $eS_0$ must
intersect.  Otherwise the right ideals $e^i dS_0$, $i\geq0$,
would form an infinite internal direct sum in $S_0$: in a finite
relation, its term of least $e$-power belongs both to $dS_0$ and
to $eS_0$ after canceling that power, and hence vanishes; continue
inductively.  This contradicts noetherianity.  The same argument
on the left gives the two-sided Ore division ring of fractions
$Q_0$.  Clearing common
denominators gives exact localization of modules.  Consequently,
for a finitely generated right module $M$, the kernel
\[
 M_{\mathrm{tor}}=\ker\bigl(M\longrightarrow M\otimes_{S_0}Q_0\bigr)
\]
consists exactly of the elements right-annihilated by a nonzero
element of $S_0$.

The quotient $M/M_{\mathrm{tor}}$ embeds in the finite-dimensional
right $Q_0$-vector space $M\otimes_{S_0}Q_0$.  Choose coordinates
in that space and finitely many generators of the embedded module.
In each coordinate, clear their common \emph{left} denominators.
Left multiplication in a coordinate is right $S_0$-linear, so this
produces an injection
\[
 M/M_{\mathrm{tor}}\longrightarrow S_0^r
\]
for some finite $r$.  The choice of left denominators is the one
compatible with right modules.

Every submodule of $S_0^r$ is free of finite rank.  To prove this,
induct on $r$ and project to the first coordinate.  The image is
a right ideal.  If it is zero, the induction hypothesis applies
inside $S_0^{r-1}$.  If it is nonzero, its principal generator is
a nonzero element of the domain, so the image is free of rank one.
The projection then splits, and its kernel lies in $S_0^{r-1}$.
This proves the assertion in both cases.  It follows that
$M/M_{\mathrm{tor}}$ is finite free, and hence that the sequence
with this quotient splits.  The kernel of a finite free presentation
of $M$ is also finite free by the same argument.  This gives a
finite free resolution of length at most one.

Now view a dg module over $S_0[t^{\pm1}]$ as a $D_0$-periodic
complex of right $S_0$-modules.  A finite-cell module has finite
free modules in each period residue, and noetherianity of $S_0$
makes its homology finitely generated there.  The same holds for
a retract, hence for every perfect module $P$.
For each $i\in\{0,\ldots,D_0-1\}$ choose a resolution
\[
 0\longrightarrow R_{1,i}\longrightarrow R_{0,i}
   \longrightarrow H_i(P)\longrightarrow0
\]
by finite free right modules.  Place $R_{0,i}$ in degree $i$ and
$R_{1,i}$ in degree $i+1$.  Lift the free generators of $R_{0,i}$
to cycles in a representative complex for $P$.  The images of
the relations are boundaries.  Since $R_{1,i}$ is free, lift them
through the differential onto the boundaries in degree $i$.
This gives a map of the two-term complex into $P$.

Extend the map periodically using multiplication by $t$.
The positive even degree $D_0$ gives the usual periodization
without changing the homological signs.  This periodized complex
has homology $H_i(P)$ in residue $i$ and zero in the other residues.
Sum the maps over $0\leq i<D_0$.  The sum is finite and induces
an isomorphism in every homology degree, so it is a
quasi-isomorphism.  Finally split each $H_i(P)$ into its torsion
and finite free parts.  The periodizations of the free parts are
finite sums of free shifts; those of the torsion parts are perfect
because of their length-one finite free resolutions.  This gives
the asserted splitting.
\end{proof}

The second puts every finite construction into a tower of free
layers, with all retracts taken at the end.

\begin{lemma}[Finite towers and retracts]
\label{rew:obs:tower-normal-form}
For an idempotent-complete triangulated category $\mathcal C$ and
an object $T_0$, every $P\in\Thick_{\mathcal C}(T_0)$ is a
retract of an object $P_0$ admitting a finite tower of triangles
\[
 F_j\longrightarrow P_j\longrightarrow P_{j+1}
      \longrightarrow F_j[1],
 \qquad 0\leq j<s,\qquad P_s=0,
\]
where each $F_j$ is a finite sum of arbitrary integer shifts of
$T_0$.
\end{lemma}

\begin{proof}
Call an object a finite extension of free layers if it admits such
a tower.  Zero, finite sums of shifts of $T_0$, and their shifts
have this property.  The octahedral axiom concatenates two towers.
More explicitly, suppose $A\to E\to B$ is a triangle, and the
first layer of a tower for $A$ is $F_0\to A\to A_1$.  The cofiber
$E_1$ of the composite $F_0\to A\to E$ fits into triangles
\[
 F_0\longrightarrow E\longrightarrow E_1,
 \qquad A_1\longrightarrow E_1\longrightarrow B.
\]
Induction on the number of layers in $A$, ending with the tower
for $B$, gives a tower for $E$.  In particular, if $A$ and $B$
have towers and $f:A\to B$ is a map, the triangle
\[
 B\longrightarrow\cofib(f)\longrightarrow A[1]
\]
gives a tower for $\cofib(f)$.  Finite sums of towers give the
same closure for finite sums.

It remains to move retracts to the end.  Let $A$ and $B$ be
retracts of finite extensions $\overline A$ and $\overline B$.
Idempotent completeness supplies complements
\[
 \overline A\simeq A\oplus A',
 \qquad \overline B\simeq B\oplus B'.
\]
Extend a map $f:A\to B$ to
$\overline f:\overline A\to\overline B$ by setting all other
matrix entries to zero.  Its cone is isomorphic to
\[
 \cofib(f)\oplus B'\oplus A'[1].
\]
Thus $\cofib(f)$ is a retract of a finite extension.  Shifts,
finite sums and successive retracts of such objects are also
retracts of finite extensions.  The resulting subcategory is thick
and contains $T_0$, which proves the lemma.
\end{proof}

\begin{proof}[Proof of Theorem~\ref{rew:obs:tower-criterion}]
Suppose that $U_0\in\Thick_{\mathcal C}(T_0)$.
By Lemma~\ref{rew:obs:tower-normal-form}, choose a retract
insertion $g_0:U_0\to P_0$ and a tower as in that lemma.
Let $g_j:U_0\to P_j$ be the composite with the map to its
$j$th tail, and put $Y_0=\Psi(U_0)$.
We prove backward from $j=s$ that
\[
 H(\Psi(g_j))=0,
 \qquad H\bigl((\Psi(g_j))^*\bigr)=0.
\]
Both assertions hold for $P_s=0$.

Write one extended step as
\[
 F\longrightarrow P_1'\longrightarrow P_2'
      \longrightarrow F[1],
 \qquad g_1':Y_0\to P_1',\quad g_2':Y_0\to P_2',
\]
where $F$ is a finite sum of free shifts, and suppose that the
two vanishing assertions hold for $g_2'$.
By Lemma~\ref{obs:periodic-splitting}, choose
\[
 P_2'\simeq V'\oplus L',
\]
with $V'$ a finite sum of free shifts and $H(L')$ torsion.
Each component of the projection $P_2'\to V'$ is a class in
the homology of $(P_2')^*$.  Dual-homology vanishing for $g_2'$
says that its composite with each of these classes is zero.
Hence
\[
 Y_0\xrightarrow{g_2'}P_2'\longrightarrow V'
\]
is zero as a derived map.  The exact sequence of maps into the
fiber triangle therefore lifts $g_1'$ to
\[
 Z'=\fib(P_1'\longrightarrow V').
\]
The fiber construction, equivalently the octahedral axiom, gives
a triangle
\[
 F\longrightarrow Z'\longrightarrow L'\longrightarrow F[1].
\]

In its homology sequence, the map preceding
$H(F)\to H(Z')$ has torsion source and free target in each
degree.  It is zero, so $H(F)\to H(Z')$ is injective.
The lifted map $Y_0\to Z'$ is zero on homology after composition
with $Z'\to L'$, because that composite is the corresponding
summand of $g_2'$.  Its homology therefore factors through
the submodule $H(F)\subset H(Z')$.  The source $H(Y_0)$ is
torsion, so its map into the free module $H(F)$ is zero.
It follows that $H(g_1')=0$.

The map $g_1'$ is the image of a map $U_0\to P_j$ in
$\mathcal C$, and this tower tail belongs to
$\Thick_{\mathcal C}(T_0)$.  Hypothesis
\eqref{rew:obs:two-tests} gives $H((g_1')^*)=0$ as well.
This completes the backward induction.  In particular
$\Psi(g_0)$ is zero on homology.  It has a left inverse because
$g_0$ is a retract insertion and every functor preserves its
retraction equation.  Thus $H(\Psi(g_0))$ is also a split
injection, contradicting the nonzero homology of $Y_0$.
\end{proof}

The two tests have distinct roles in the induction.  Dual homology
removes the free quotient of a tower tail as an actual derived map.
Ordinary homology then places the remaining map in a free submodule,
where torsion forces it to vanish.  The implication
\eqref{rew:obs:two-tests} restores the dual statement for the next
step.  There is no bound on the length of the proposed tower.

\subsection{A flat coefficient extension}

We now construct the coefficient ring required by the criterion.
The extension must be flat on both sides: homology of right
modules and homology of their left-module duals will both be tested
after extending coefficients.  It must also keep the cone of $a$
nonzero while making its homology torsion.

There are two filtrations in what follows.  The paired vectors in
$W_i(Y)=\mathcal Q\oplus\mathcal K$ have Laurent pole bounds in
$K_0=k((x))$.  The operation ring $\Delta=R[[H]]$, with $R=k[[x]]$,
has instead the weight filtration of
Proposition~\ref{geom:skew-filtration}, with
\[
 \gr\Delta=k[\sigma x,A,B,C],\qquad
 |\sigma x|=4,\quad |A|=|B|=1,\quad |C|=3.
\]
The coefficient extension will be taken at the homogeneous prime
$(C)$ and used only to test homology of perfect modules.  The later
estimate concerns integral operators acting on
Laurent functions: it compares the growth of their joint kernels
with the larger exponent $\beta>1/3$ supplied by
Lemma~\ref{glue:pairs}.

We give the filtered algebra argument with explicit hypotheses.  A
decreasing filtration is indexed by integers and is multiplicative.  Its
initial-symbol map is denoted by $\sigma$.  A filtration on a finitely
generated module is called \emph{good} if it is the quotient filtration
from a finite direct sum of filtration shifts of the ring.
The construction through truncated Rees rings is the algebraic approach
to microlocalization developed by Asensio, Van den Bergh, and
Van Oystaeyen~\cite[Sections~2--3]{AVVO1989}.  We give the construction
and flatness argument in the precise complete setting needed here.

\begin{lemma}[Filtered microlocalization]
\label{obs:microlocal-flat}
Let $k$ be a finite field, and let $\mathcal A$ be a complete separated
filtered $k$-algebra such that $F^n\mathcal A=\mathcal A$ for $n\leq0$
and
\[
 \gr\mathcal A=k[z_1,\ldots,z_d,c]
\]
as a graded algebra, with all the displayed generators of positive
integer degree.  There is a filtered algebra $\mathcal S$ and a map
$\mathcal A\to\mathcal S$ with the following properties:
\begin{enumerate}
 \item $\gr\mathcal S$ is the homogeneous localization of
 $\gr\mathcal A$ obtained by inverting all homogeneous elements outside
 $(c)$;
 \item $\mathcal S$ is separated and complete for its exhaustive
 filtration, which is allowed to have nonzero terms of negative degree;
 \item $\mathcal S$ is flat as both a left and a right
 $\mathcal A$-module.
\end{enumerate}
For a finitely generated good filtered right module $M$, the corresponding
microlocalized module is $M\otimes_{\mathcal A}\mathcal S$ and its
associated graded is the same homogeneous localization of $\gr M$.
The analogous assertion holds for left modules.
\end{lemma}

\begin{proof}
We first justify the finite-module filtration properties needed for the
construction.  If $N$ is a submodule of a finite shifted filtered free
module, choose homogeneous generators of $\gr N$ over the noetherian
ring $\gr\mathcal A$ and lift them to $N$.  For an element of $N$,
subtract a linear combination of the lifts matching its initial symbol,
and repeat.  The remainder increases in filtration at each step.
The coefficients of each chosen lift form a convergent series in
$\mathcal A$, so the element is a finite linear combination of the
lifts.  This also proves that the resulting surjection from a finite
shifted free module is strict for the induced filtration.

Each quotient $\mathcal A/F^n\mathcal A$ is finite: the graded pieces
below $n$ are finite-dimensional over the finite field $k$.
Consequently $\mathcal A$ and its finite products are compact.
The submodule just considered is therefore compact and closed, being
the image of a compact finite free module.  Notice that closedness was
not assumed during the initial-symbol subtraction.  These observations
prove that $\mathcal A$ is noetherian on both sides and that every
finitely generated module has a complete separated good filtration.
Submodule and quotient filtrations are good, by using their preimages
in a finite free presentation.  A map of finite good filtered modules
is bounded up to a filtration shift, as one sees on strict generators.
In particular, any two good filtrations on one module differ by bounded
shifts.

Form the Rees ring and Rees module
\[
 \widetilde{\mathcal A}
   =\bigoplus_{n\in\mathbb Z}F^n\mathcal A\,s^{-n},
 \qquad
 \widetilde M
   =\bigoplus_{n\in\mathbb Z}F^nM\,s^{-n},
\]
where $s$ is central of Rees degree $-1$.
For $j\geq1$, work modulo $s^j$ and invert the homogeneous elements
whose reductions modulo $s$ lie outside $(c)$.
Such elements are regular, by induction through the $s$-layers and
regularity of their nonzero reductions in the polynomial ring.
They satisfy the Ore conditions on both sides.  Indeed all commutators
are divisible by $s$, so for a denominator $d$ and an element $e$ one
has $\operatorname{ad}_d^j(e)=0$ modulo $s^j$.  Expanding this identity
shows that $d^j e$ has a right factor $d$, and that $e d^j$ has a left
factor $d$.  Powers of a denominator are still denominators.  These
relations give the Ore conditions first for homogeneous elements and
then for their finite sums.  We may thus use ordinary two-sided Ore
localization, retaining the Rees grading.  Localization is exact on
modules, by clearing common denominators.

Let $M_j^\mu$ denote this localization of
$\widetilde M/s^j\widetilde M$.  Denominators lift across truncation
levels.  The localizations therefore reduce compatibly, and localizing
a module from a lower level at a higher level gives its own
localization.  Define
\begin{equation}
\label{obs:mu-definition}
 F^n\mu(M)=\lim_j(M_j^\mu)_n,
 \qquad
 \mu(M)=\bigcup_{n\in\mathbb Z}F^n\mu(M).
\end{equation}
The inclusions in this union come from multiplication by powers of $s$.
More precisely, regularity of $s$ on the untruncated Rees module gives,
after exact localization,
\begin{equation}
\label{obs:rees-layers}
 0\longrightarrow(M_j^\mu)_{n+l}
 \xrightarrow{\ s^l\ }(M_{j+l}^\mu)_n
 \longrightarrow(M_l^\mu)_n\longrightarrow0
 \qquad(l,j\geq1).
\end{equation}
Transitions in each fixed Rees degree are surjective.  Taking limits
in \eqref{obs:rees-layers} gives the asserted inclusions and
\[
 F^n\mu(M)/F^{n+l}\mu(M)=(M_l^\mu)_n.
\]
The filtration is separated and complete.  In fact compatible Cauchy
data eventually lie in one fixed $F^n\mu(M)$, where both assertions
follow from the inverse-limit description.  Taking $l=1$ identifies
the associated graded with the homogeneous localization of $\gr M$.

Products at the truncated levels give
$\mathcal S=\mu(\mathcal A)$ its ring structure and give $\mu(M)$
its module structure.  The natural maps from $\mathcal A$ and $M$
are obtained by taking their classes at each level.  A bounded-shift
map of good filtered modules gives a map on these unions; changing
the bound does not change the map, by the transition inclusions.
Thus the underlying module $\mu(M)$ is functorial and independent
of the good filtration chosen.

This functor is exact on finite modules.  Use induced and quotient
filtrations in a short exact sequence to make it strict.  Its Rees
sequence remains exact modulo $s^j$, because the quotient Rees module
is $s$-torsion-free.  Ore localization preserves exactness.  Inverse
limits preserve it here because the transition maps in each module
system are surjective, and taking the filtered union preserves it as
well.  The natural map
\[
 M\otimes_{\mathcal A}\mathcal S\longrightarrow\mu(M)
\]
is an isomorphism for a free module and hence for every finite module
by a finite presentation.  This proves left flatness of $\mathcal S$;
over a noetherian ring, exactness on finite modules suffices for this
flatness test.  The identical construction for left modules proves
right flatness.
\end{proof}

Apply Lemma~\ref{obs:microlocal-flat} to $\mathcal A=\Delta$ and
$c=C$, and write $S$ for the resulting ring.  The initial symbol of
$x$ is a unit in $\gr S$.  Lifting its inverse and correcting the
successive errors in increasing filtration shows that $x$ is a unit
in $S$.  Both left and right flatness hold also over $\Delta[x^{-1}]$:
on a module over that ring, balancing by the invertible $x$ identifies
tensoring over $\Delta$ with tensoring over $\Delta[x^{-1}]$.

Recall from Proposition~\ref{complete:operations} that
\[
 \mathscr E[t^{-1}]=A_{\mathrm{op}}[t^{\pm1}],
 \qquad A_{\mathrm{op}}=(\Delta[x^{-1}])^\Lambda.
\]
The tame character decomposition makes $\Delta[x^{-1}]$ finite free
over $A_{\mathrm{op}}$ on either side.  Indeed $x$ has a faithful
$\Lambda$-character and is a unit; multiplication on either side
by its powers identifies every character space with the invariant
space.  We have proved the following coefficient-extension statement:
\begin{equation}
\label{obs:coefficient-extension}
 \mathscr E\longrightarrow
 A_{\mathrm{op}}[t^{\pm1}]
 \longrightarrow\Delta[x^{-1}][t^{\pm1}]
 \longrightarrow\mathscr B:=S[t^{\pm1}]
\end{equation}
is flat on both sides.  Here $t$ is central, of homological degree $D$;
$S$ is in homological degree zero and has zero differential.
The filtration on $S$ is unrelated to this homological grading.

The following lemma verifies the remaining ring properties and
shows that the cone of $a$ retains nonzero torsion homology.

\begin{lemma}
\label{obs:principal}
The ring $S$ is a domain, every left or right ideal of $S$ is principal,
and $S$ has a two-sided Ore division ring of fractions.  The operator
$a$ constructed in Proposition~\ref{op:operator} has no one-sided inverse
in $S$.  In particular $S/aS\ne0$, and this right $S$-module is torsion.
\end{lemma}

\begin{proof}
The graded ring of $S$ is a domain, so multiplicativity of initial
symbols proves that $S$ is a domain.  Every nonzero homogeneous element
of $\gr S$ is a power of $C$ times a homogeneous unit.  In a nonzero
right ideal choose $d_0$ minimizing the nonnegative $C$-order of its
initial symbol.  For any member $d$ of the ideal, its initial symbol
is a multiple of $\sigma(d_0)$.  Lift the multiplier and subtract the
corresponding right multiple of $d_0$.  The remainder is still in the
ideal and has strictly larger filtration, so its initial $C$-order
is still at least the chosen minimum.  Repeating, and summing the
convergent multipliers, expresses $d$ as a right multiple of $d_0$.
The successive multipliers tend to zero because their filtration
degrees increase by the same amount as those of the remainders.
This proves the right principal-ideal property; the left proof is
identical.

The intersection argument in the proof of
Lemma~\ref{obs:periodic-splitting} now gives the two-sided Ore
division ring of fractions.

The operator $a$ still has initial symbol $C$ in $S$.  A one-sided
inverse would make $C$ invertible in $\gr S$, which it is not.
Thus $S/aS\ne0$.  Given a nonzero representative $b\in S$, the
nonzero intersection $bS\cap aS$ supplies a nonzero $d$ with
$bd\in aS$.  Every class in $S/aS$ is therefore annihilated on the
right by a nonzero element, proving torsion.
\end{proof}

\paragraph{The instance supplied by the glued object.}
We now return to $\mathbf D(\mathcal N)$ and the coefficient
extension \eqref{obs:coefficient-extension}.  Let
\[
 \mathcal C_{\mathrm{comp}}
 =\{P\in\mathbf D(\mathcal N):
       \widehat P\in\Thick(\widehat T')\}
\]
be the full subcategory with the displayed property.
Completion is exact by Proposition~\ref{complete:fracture}, so
the inverse image of $\Thick(\widehat T')$ is thick.  The ambient
derived category is idempotent-complete, and the defining property
is closed under retracts; hence $\mathcal C_{\mathrm{comp}}$ is
idempotent-complete as well.  It contains $T'$, and
it contains $U$ because Proposition~\ref{glue:object} identifies
$\widehat U$ with the displayed finite construction $Y$.
Compose completion, the inverse of the strict realization
equivalence in Proposition~\ref{complete:realization}, and derived
extension of scalars to obtain an exact functor
\[
 \Psi:\mathcal C_{\mathrm{comp}}\longrightarrow\Perf(\mathscr B).
\]
We write $\widehat P_{\mathscr B}$ for $\Psi(P)$ and use the
same subscript for the image of a map.  The functor sends $T'$
to the free right module $\mathscr B$.

On a free right module, the map determined by $1\mapsto a$ is
left multiplication by $a$.  It is injective over $S$ because
$S$ is a domain, and its cokernel is the right module $S/aS$.
Since $Y$ is the sum of all $D$ shifts of this cone,
\begin{equation}
\label{rew:obs:source-homology}
 H_i(\Psi(U))\cong S/aS
 \quad\text{as right $S$-modules for every integer }i.
\end{equation}
Lemma~\ref{obs:principal} says that this module is nonzero and
torsion.  This verifies the first hypothesis of
Theorem~\ref{rew:obs:tower-criterion}.  The remaining task is
the implication from homology vanishing to dual-homology
vanishing for every global map $g:U\to P$ with
$P\in\Thick_{\mathbf D(\mathcal N)}(T')$.

\subsection{A growth bound for joint operator kernels}

To prove the remaining implication, we return to integral operators
acting on $K_0=k((x))$.  For the remainder of this section let
$\bar b$ denote the reduction of $b\in\Delta$ in
$\Omega=\Delta/(x)=k[[H]]$.  The notation $\sigma(\bar b)$ means
the subsequent initial symbol for the weight filtration on
$\Omega$; reduction modulo $x$ and passage to the initial symbol
are distinct operations.
The congruence subgroups $H_l$, their indices $N_l$, and the balanced
sequence $b_j^H$ are those used in
Lemma~\ref{op:displacement} and
Lemma~\ref{op:divided-differences}.  In particular, with
$z_h=h(x)/x$, we have
\begin{equation}
\label{obs:growth-inputs}
 r'_m=\sum_{l\geq1}\lfloor m/N_l\rfloor\leq\kappa m,
 \qquad
 \Gamma_s\subset H_{p^s},\qquad
 N_{p^s}\geq p^{3(s-1)}\quad(s\geq1).
\end{equation}
These are the inputs that will turn linear Laurent precision into
cubic-root dimension growth.

\begin{lemma}[Joint-kernel growth]
\label{obs:growth}
Suppose $a_1,b_1\in\Delta$ satisfy
\[
 \bigl(\sigma(\bar a_1)\bigr)=(C)
 \quad\text{in }\gr\Omega=k[A,B,C],
 \qquad \sigma(\bar b_1)\notin(C).
\]
Fix an integer $L_0\geq0$.  There is a constant $c_0$, depending on
the fixed chart, $L_0$, $a_1$, and $b_1$, such that for every $n\geq1$
the image of
\[
 \{f\in x^{-L_0}R:a_1(f)=b_1(f)=0\}
       \longrightarrow K_0/x^nR
\]
has $k$-dimension at most $c_0(1+n)^{1/3}$.
\end{lemma}

\begin{proof}
Choose any finite family $f_1,\ldots,f_v$ in the indicated joint
kernel which is independent modulo $x^nR$.  We bound $v$ uniformly.
Consider the orbit functions
\[
 \varphi_i(h)=h(f_i),\qquad h\in H,
\]
as functions with values in $K_0$.  The proof first places them
in a lattice of integral-valued functions whose reduction injects
into functions on $H$.  Interpolation then makes that reduction
factor through a finite quotient, where the two annihilator
relations give the required dimension bound.

We first bound the coefficients of integral-valued linear
combinations of the $\varphi_i$.  Put
$F_i(y)=y^{L_0}f_i(y)\in k[[y]]$.
For coefficients $c_i\in K_0$ with minimum $x$-valuation zero,
\[
 \sum_i c_i \varphi_i(h)
   =h(x)^{-L_0}\sum_i c_i F_i(h(x)).
\]
The reduction modulo $x$ of $\sum_i c_iF_i(y)\in R[[y]]$ is a
nonzero series of order $d<n+L_0$.  Indeed the residues of the $c_i$
give a nontrivial $k$-linear combination, and the assumed independence
of the $f_i$ bounds its order after multiplication by $y^{L_0}$.
Weierstrass preparation writes this series as a unit times a monic
polynomial of degree $d$.  Evaluation at $h(x)=xz_h$ preserves
valuation of the polynomial because the evaluated unit is a unit.

As a polynomial in $z_h$, that monic polynomial evaluated at $xz_h$
has leading coefficient $x^d$.  Lagrange interpolation on the first
$d+1$ balanced orbit points expresses $x^d$ as the sum of its values
divided by the interpolation denominators.  Every such denominator
has valuation at most $r'_d$.  Therefore at least one of the values
has valuation at most $d+r'_d$.  By
\eqref{obs:growth-inputs}, some value of the normalized orbit
combination has valuation $O(n)$, with constants independent of
the chosen family and coefficients.

It follows first that the $\varphi_i$ are independent over $K_0$.
It follows also that the coefficients of every integral-valued
combination have valuations bounded below by $-O(n)$: factor out
the smallest coefficient valuation and apply the preceding upper
bound to the remaining normalized combination.
Let $\mathcal L_{\mathrm{orb}}$ be the $R$-module of integral-valued functions in
the $K_0$-span of the $\varphi_i$.  It contains each $x^{L_0}\varphi_i$ and is
contained in a fixed finite free lattice by the coefficient bound.
Thus it is a full finite free lattice of rank $v$.
Its reduction injects into the space of functions $H\to k$.
In fact an integral-valued function whose values all lie in $xR$
is divisible by $x$ within $\mathcal L_{\mathrm{orb}}$, by its definition.

Multiply an integral combination by $z_h^{L_0}$, which is congruent
to one modulo $x$.  Its residue is consequently the residue of
\[
 x^{-L_0}\sum_i c_iF_i(xz_h).
\]
The uniform coefficient bound permits truncation to a polynomial
in $z_h$ of degree $J=O(n)$: all later terms vanish modulo $x$.
The truncated polynomial is integral-valued on the whole orbit,
since it differs from the preceding integral expression by an
$xR$-valued function.

We next control the congruence level on which the residue of such
a polynomial depends.  The normalized Newton basis is the
integral-valued-polynomial construction associated with $P$-orderings;
see \cite[\S11, Theorem~24]{Bhargava2000} for the general framework.
The orbit-specific bounds and residue estimates below are proved here.
The integral-valued polynomials of degree
at most $J$ on the orbit points have the $R$-basis
\begin{equation}
\label{obs:newton-basis}
 x^{-r'_m}\prod_{i=0}^{m-1}(z-z_{b_i^H}),
 \qquad 0\leq m\leq J.
\end{equation}
Balancedness gives at least $\lfloor m/N_l\rfloor$ factors
congruent to zero modulo $x^l$ at every orbit point.  Hence each
polynomial in \eqref{obs:newton-basis} is integral-valued.
At $z=z_{b_m^H}$ these counts are exact, so its value is a unit.
Successively evaluating at $z_{b_0^H},z_{b_1^H},\ldots$ now proves
the basis assertion by triangular interpolation.

If two orbit points are congruent modulo $x^l$ and $N_l>J$, the
residues of all the polynomials in \eqref{obs:newton-basis} agree
at the two points.  To see this, expand the difference of a product
by replacing one factor at a time.  Deleting that factor loses
at most one of the required matches at each level below $l$.
No matches at levels $l$ and higher are needed, since $N_l>m$.
The difference of the replaced factors contributes valuation at
least $l$, exceeding the possible loss of $l-1$.  After division
by $x^{r'_m}$ every difference is thus in $xR$.

Choose $s\geq1$ with $N_{p^s}>J$ and
$p^s=O(n^{1/3})$.  Such a choice follows from the last inequality
of \eqref{obs:growth-inputs}, by taking the first $s$ for which
$p^{3(s-1)}>J$.  All the residue functions coming from
$\mathcal L_{\mathrm{orb}}$ then factor through $H/\Gamma_s$.

We have placed the rank-$v$ lattice, after reduction, inside
functions on a finite quotient.  The annihilator equations now
impose two algebraic relations on that function space.
For any $\varphi\in\mathcal L_{\mathrm{orb}}$ and $\mu=a_1,b_1$, the original
annihilator equations imply
\begin{equation}
\label{obs:orbit-equations}
 \int_{g\in H}h\bigl(\mu(g)\bigr)\varphi(hg)=0
 \qquad(h\in H).
\end{equation}
For the functions $\varphi_i$ this is obtained by applying $h$ to
$\mu(f_i)=0$; field-linear combinations preserve the identity
because their coefficients do not depend on $h$.
All integrals can be computed at finite coefficient and group
precision.  On reduction, $h$ acts trivially on the coefficient
field $k$.  The resulting functions on $H/\Gamma_s$ therefore
annihilate the left translates of $\bar a_1$ and $\bar b_1$.
Their space has dimension at most
\[
 \dim_k\frac{\Omega}
 {\Omega\bar a_1+\Omega\bar b_1+
        \ker\bigl(\Omega\longrightarrow k[H/\Gamma_s]\bigr)}.
\]
The associated graded of this quotient kills at least
\[
 C,\quad \sigma(\bar b_1),\quad
 A^{p^s},\quad B^{p^s},\quad C^{p^s},
\]
since $g_i^{p^s}\in\Gamma_s$.  The quotient filtration is finite,
already on the finite group algebra.  Modulo $C$, take a leading
monomial $A^{u_0}B^{v_0}$ of the fixed nonzero polynomial
$\sigma(\bar b_1)\bmod C$.  Among the monomials
$A^iB^j$ with $0\leq i,j<p^s$, at most $(u_0+v_0)p^s$ are not
divisible by $A^{u_0}B^{v_0}$; the case of a nonzero constant gives a zero
quotient.  Thus the displayed dimension is $O(p^s)=O(n^{1/3})$.
Because $\mathcal L_{\mathrm{orb}}/x\mathcal L_{\mathrm{orb}}$ has dimension $v$ and injects
into this function space, the required bound follows.
\end{proof}

We now have a quantitative restriction on joint kernels of two
integral operators.  To apply it to a global map out of $U$, we must
extract a second annihilating operator from a microlocal homology
vanishing condition.  The next lemma records both sides of this
extraction; they will be used separately and must not be interchanged.

\begin{lemma}[Removal of $x$-divisible remainders]
\label{obs:right-annihilator}
Let $a\in\Delta$ be the operator of
Proposition~\ref{op:operator}, so that $\bar a=U_3$ and
$\sigma(\bar a)=C$.
\begin{enumerate}
 \item Suppose a right ideal $J\subset\Delta$ contains $a\Delta$,
 is right $x$-saturated, meaning $bx\in J\Rightarrow b\in J$,
 and satisfies $JS=S$.  Then $J$ contains an element $b$ such that
 $\sigma(\bar b)\notin(C)$.
 \item\label{obs:left-reduction}
 Suppose $b_0\in\Delta[x^{-1}]$ does not belong to $Sa$ and
 annihilates a $k$-subspace $V\subset\ker(a:K_0\to K_0)$.
 There is $b_1\in\Delta$ annihilating $V$ with
 $\sigma(\bar b_1)\notin(C)$.
\end{enumerate}
\end{lemma}

\begin{proof}
For the first assertion, suppose that every member of $J$ reduces
into $\bar a\Omega$.  Starting with $j\in J$, subtract a right
multiple of $a$ with the same reduction and write
\[
 j=a d_0+j_1x.
\]
Normality of $x$ allows the right division, and right saturation
puts $j_1$ back in $J$.  Repetition gives
\[
 j=a(d_0+d_1x+\cdots+d_{r-1}x^{r-1})+j_rx^r.
\]
The remainder tends to zero, and the coefficient series converges
in $\Delta$.  Thus $j\in a\Delta$, and $J=a\Delta$.
This contradicts $JS=S$, because $aS\ne S$ by
Lemma~\ref{obs:principal}.

Choose now $j\in J$ whose reduction is not in $\bar a\Omega$.
If its reduction has initial symbol divisible by $C$, subtract a
right multiple of $\bar a$ matching that symbol.  Continue in the
complete filtered ring $\Omega$.  The process must terminate with
an initial symbol outside $(C)$: otherwise the convergent sum of
multipliers would express the original reduction as a member of
$\bar a\Omega$.  Lift the finitely many subtractions to right
multiples of $a$ in $\Delta$.  They preserve membership in $J$
and produce the asserted $b$.

For the second assertion, first multiply $b_0$ on the left by a
power of $x$ to obtain $b\in\Delta$.  This preserves annihilation
of $V$ and nonmembership in $Sa$, since $x$ is a unit on $K_0$
and in $S$.  Whenever $\bar b\in\Omega\bar a$, subtract a left
multiple of $a$ and divide on the left by $x$.  If this procedure
never terminated, its equations would give
\[
 b=(d_0+xd_1+\cdots+x^{r-1}d_{r-1})a+x^r b_r,
 \qquad b_r,d_r\in\Delta.
\]
The limit would put $b$ in $\Delta a$, a contradiction.  Thus a
reduction outside $\Omega\bar a$ is reached.  As on the right,
successive initial-symbol subtraction by left multiples of
$\bar a$ then terminates with a symbol outside $(C)$.
Throughout, subtracting a left multiple of $a$ preserves
annihilation of $V$, and division on the left by $x$ does also.
The final integral operator is the required $b_1$.
Neither of these two convergent expansions commutes $x$ through $a$.
\end{proof}

\subsection{Dual detection for global maps}

We now verify the second hypothesis of
Theorem~\ref{rew:obs:tower-criterion} for the functor
$\Psi(P)=\widehat P_{\mathscr B}$ defined above.
Proposition~\ref{complete:realization} identifies each complete
object in $\Thick(\widehat T')$ with a perfect right
$\mathscr E$-module.  For such a module $M$, the two duals used
below are
\[
 M^*=R\Hom_{\mathscr E}(M,\mathscr E),
 \qquad
 (M_{\mathscr B})^*
   =R\Hom_{\mathscr B}(M_{\mathscr B},\mathscr B).
\]
They are left modules.  The subscript $\mathscr B$ denotes the
extension \eqref{obs:coefficient-extension}.
A \emph{homology ghost} here is a map inducing zero on module
homology in every degree.  We will show that a global map out of
$U$ which becomes a homology ghost also becomes zero on dual
homology.

\begin{lemma}[Adjoints on the Laurent field]
\label{obs:adjoint}
Every element $b\in\Delta$ has an adjoint $b^\dagger\in\Delta$
for the continuous pairing
\[
 \langle u,f\rangle=\CT_x(uf),\qquad u,f\in K_0.
\]
Here $\CT_x$ extracts the coefficient of $x^0$ in a Laurent series.
Modulo $x$, taking the adjoint induces group inversion on
$\Omega=k[[H]]$.  In particular it preserves the conditions
that an initial symbol generate $(C)$ or lie outside $(C)$.
\end{lemma}

\begin{proof}
The pairing identifies $K_0$ with its continuous $k$-linear dual.
Indeed a continuous functional kills some $x^jR$; its values on
the monomials $x^i$ are exactly the coefficients of a Laurent
series bounded below, read in the reverse order.
For a coefficient $c_0\in R$ and a group element $g\in H$,
formal residue substitution in $dx/x$ gives
\begin{equation}
\label{obs:adjoint-formula}
 (c_0\delta_g)^\dagger
 =g^{-1}(c_0)\,
   \frac{x\,(g^{-1}(x))'}{g^{-1}(x)}\,
   \delta_{g^{-1}}.
\end{equation}
The Jacobian factor is an integral unit congruent to one modulo
$x$, since $g(x)/x\in1+xR$.
Residue substitution remains valid in characteristic $p$:
for monomials, the requisite identities depend on only finitely
many coefficients of an invertible formal series and hold
integrally, as is seen by verifying them first over characteristic
zero.  They then specialize to $k$ and extend to Laurent series.

Formula~\eqref{obs:adjoint-formula} extends to measures in
$\Delta$.  At any fixed coefficient precision, moving the
coefficients by $g^{-1}$ and multiplying by the Jacobian require
only a finite level of the group action; inversion is continuous
in the group coordinates.  Equivalently, integrate the displayed
formula on sufficiently fine finite quotients, at each precision.
The resulting compatible measures define $b^\dagger\in\Delta$
and satisfy the adjoint identity.  Modulo $x$, this is group
inversion on the completed group algebra.  Its associated graded
map sends the generator symbols $A,B,C$ to their negatives.
It therefore preserves the two asserted symbol conditions.
\end{proof}

\begin{proposition}[Dual detection of global ghosts]
\label{obs:dual-ghost}
For the object $U$ of Proposition~\ref{glue:object}, let
$g:U\to P$ be a map in $\mathbf D(\mathcal N)$ with
$P\in\Thick(T')$.  Then
\begin{equation}
\label{obs:ghost-implication}
 H(\widehat g_{\mathscr B})=0
 \quad\Longrightarrow\quad
 H\bigl((\widehat g_{\mathscr B})^*\bigr)=0.
\end{equation}
Both zero assertions mean that the corresponding induced homology
maps vanish.
\end{proposition}

\begin{proof}
Assume the left side of \eqref{obs:ghost-implication}.
We first bound the images of the polar representatives belonging
to the cokernel part of $W(Y)$, then use the paired lattice to
obtain large killed subspaces of the paired span.  Composites into
the generator vanish on the cokernel part, so they annihilate
the kernel projections of these subspaces.  Lemma~\ref{obs:growth}
will force the operators representing those composites into the
left ideal $Sa$.

\paragraph{An integral second annihilator.}
Fix a residue $r$ and let
$\iota_r:\widehat T'[r]\to m[r]$ be the cofiber inclusion.
Restrict $\widehat g$ to this summand of $Y$ and put
\[
 h:\widehat T'[r]\xrightarrow{\iota_r}m[r]
       \longrightarrow Y\xrightarrow{\widehat g}\widehat P.
\]
This map is a homology class of the perfect module representing
$\widehat P$.  Extend coefficients first to
$\Delta[x^{-1}][t^{\pm1}]$, but not yet to $\mathscr B$, and
let $J_h\subset\Delta$ be the right annihilator of that class.
The ring $\Delta$ acts in homological degree zero here.
We have $a\Delta\subset J_h$ because $\iota_r a=0$.
The ideal is right $x$-saturated because $x$ acts invertibly on
the extended module.

Put $A_{\mathrm{loc}}=\Delta[x^{-1}]$ and
$N_h=H_r(\widehat P_{A_{\mathrm{loc}}[t^{\pm1}]})$, and let
$I_h=\operatorname{Ann}_{A_{\mathrm{loc}}}(h)$ be the local right
annihilator.  Then $J_h=I_h\cap\Delta$ and
$I_h=J_hA_{\mathrm{loc}}$: every local coefficient is a right
fraction $d x^{-n}$, and $h d x^{-n}=0$ if and only if $h d=0$,
by multiplying on the right by the unit $x^n$.
The cyclic map
\[
 A_{\mathrm{loc}}/I_h\longrightarrow N_h,\qquad
 \alpha+I_h\longmapsto h\alpha
\]
is injective.  Left flatness of $S$ over $A_{\mathrm{loc}}$
preserves this injection, with source
$(A_{\mathrm{loc}}/I_h)\otimes_{A_{\mathrm{loc}}}S\cong S/I_hS$.
The same left flatness, after adjoining $t$, commutes extension
with homology.  It gives
\[
 0\longrightarrow S/I_hS
 \longrightarrow N_h\otimes_{A_{\mathrm{loc}}}S
 \cong H_r(\widehat P_{\mathscr B}).
\]
The generator maps to the extended class $h_{\mathscr B}$, which
is zero because it factors through the homology ghost
$\widehat g_{\mathscr B}$.  Injectivity forces $1+I_hS=0$, so
$J_hS=I_hS=S$.
Lemma~\ref{obs:right-annihilator} gives
\begin{equation}
\label{obs:second-annihilator}
 b\in J_h,\qquad \sigma(\bar b)\notin(C).
\end{equation}
Crucially, $b$ annihilates the homology class \emph{before}
microlocal extension; $S$ was used to prove its existence, not
to define a new action on functions.

\paragraph{From the annihilator to a Laurent rank bound.}
Extend the overlap calculation from $\Phi$ to $K_0$ by derived
tensor with the left module $K_0[t^{\pm1}]$ over
$\Delta[x^{-1}][t^{\pm1}]$.  The integral skew action extends to
Laurent functions by the same integration formula: for a fixed
bounded input lattice and a fixed output precision, finite
coefficient and group levels suffice.
The relations $ha=hb=0$ are zero homology classes of maps out of
a free shift, hence zero derived maps before tensoring.  Derived
tensor therefore preserves them, without any flatness assumption
on $K_0$ as a coefficient module.

To make the return to the original overlap explicit, put
$D_{\rm loc}=\Delta[x^{-1}][t^{\pm1}]$ and let $\mathsf P$
be the perfect right $\mathscr E$-module representing $\widehat P$.
Derived associativity gives the natural identification
\[
 (\mathsf P\otimes^{\mathbf L}_{\mathscr E}D_{\rm loc})
       \otimes^{\mathbf L}_{D_{\rm loc}}K_0[t^{\pm1}]
 \simeq
 \mathsf P\otimes^{\mathbf L}_{\mathscr E}K_0[t^{\pm1}].
\]
The same identification applies to the source free shift
$\mathscr E[r]$ and is natural for $h$.
The tame average
\[
 \operatorname{pr}:K_0\longrightarrow K_0^\Lambda=\Phi,
 \qquad f\longmapsto\frac1{|\Lambda|}
                  \sum_{\lambda\in\Lambda}\lambda(f)
\]
is $\mathscr E$-linear, because the normalized operations in
$\mathscr E$ are $\Lambda$-invariant. Tensoring this map over
$\mathscr E$ projects the last displayed complex onto
$\mathsf P\otimes^{\mathbf L}_{\mathscr E}\Phi[t^{\pm1}]$,
and similarly for the source. Naturality gives a commuting square
for these two projections and the maps induced by $h$.
Applying this projection \emph{after} the preceding derived
vanishing yields
\begin{equation}
\label{obs:projected-annihilation}
 W(h)\bigl(\operatorname{pr}(a(f))\bigr)
 =W(h)\bigl(\operatorname{pr}(b(f))\bigr)=0
 \qquad(f\in K_0).
\end{equation}
The operator $b$ need not be invariant, and $\operatorname{pr}$ need
not be linear over $\Delta[x^{-1}]$.  The extended tensor complex
is equally the original perfect $\mathscr E$-module tensored
with $K_0[t^{\pm1}]$, so the $\mathscr E$-linear projection is
all that is required.

Choose a fixed compact open subspace $C_P$ in the target overlap
degree.  Continuous functionals on that overlap which kill $C_P$
pull back along $W(h)\operatorname{pr}$ to functionals on $K_0$.
By Lemma~\ref{obs:adjoint}, write them as $\CT_x(uf)$.
The map $W(h)$ is continuous by
Proposition~\ref{complete:topology}: both $T'$ and $P$ have
degreewise finite sections homotopy, and the reduction maps give
the intrinsic topology on their completions. Since the tame average
is also continuous, $W(h)\operatorname{pr}$ sends some $x^jR$
into $C_P$. All the pulled-back functionals therefore kill
$x^jR$. In the constant-term pairing this means
$u\in x^{1-j}R$, so all these $u$ lie in one fixed
$x^{-L_0}R$, with $L_0=\max(0,j-1)$.
Equation~\eqref{obs:projected-annihilation} and nondegeneracy
of the pairing show that
\[
 a^\dagger(u)=b^\dagger(u)=0.
\]
The adjoint lemma and \eqref{obs:second-annihilator} verify the
symbol hypotheses of Lemma~\ref{obs:growth}.

It follows that, modulo $C_P$, the images under $W(h)$ of the
selected polar representatives $w^j$, $1\leq j\leq N$, have
rank $O(N^{1/3})$.  Indeed take independent functionals on their
span in the discrete quotient and extend them to that quotient.
Their pullback series $u$ are independent modulo $x^{qN+1}R$:
any linear combination lying in this subspace pairs to zero with
every $w^j=x^{-qj}$, $j\leq N$, contradicting independence on
the selected images.  Lemma~\ref{obs:growth}, with
$n=qN+1$, gives the rank bound.  This argument applies in every
period residue.

\paragraph{The paired lattice forces a large killed subspace.}
In a fixed degree identify
$W_i(Y)=\mathcal Q\oplus\mathcal K$ as in
Lemma~\ref{glue:overlap}.  The images of $(q_j,0)$ are
the images just bounded, through the relevant map $W(h)$.
All images of $(0,k_l)$ belong to one compact open subspace of
the target.  In fact the $k_l$ lie in the compact subspace
$\mathcal K\cap R$ of the source, the map is continuous, and
any compact subset of the target lies in a sufficiently large
compact open stage.  Enlarge $C_P$ to contain those images.

By Lemma~\ref{glue:global-maps}, the sums $(q_j,k_l)$ belong
after applying $W(\widehat g)$ to the discrete lattice image of
$\pi_iLP$.  At scale $N$, Lemma~\ref{glue:pairs} supplies at
least $cN^\beta$ simultaneous pairs.  The rank of their images
modulo $C_P$ is $O(N^{1/3})$, while the intersection of the
discrete lattice with $C_P$ has fixed finite dimension.
Therefore the span of these pairs contains a subspace killed by
$W(\widehat g)$ of dimension at least
\begin{equation}
\label{obs:killed-pair-bound}
 cN^\beta-O(N^{1/3}),\qquad c>0.
\end{equation}
Its projection to $\mathcal K$ has the same dimension even
modulo $x^{qN+1}R$, since the selected $k_l$ are independent
modulo that subspace.  These projected vectors lie in
$\mathcal K\cap R$.

\paragraph{Composites into a free shift.}
Take any map of perfect modules, before microlocal extension,
\[
 \theta:\widehat P\longrightarrow\widehat T'[\ell].
\]
On the summand $m[r]$, the composite $\theta\widehat g$ kills
the cofiber inclusion $\iota_r$.  Indeed its composite $c$ on
that free shift satisfies $ca=0$, and right multiplication by
$a$ on homogeneous operations is injective: the operation ring
embeds in the domain $\Delta[x^{-1}][t^{\pm1}]$.
The cofiber triangle therefore factors this composite as
\[
 m[r]\xrightarrow{\partial_r}\widehat T'[r+1]
             \xrightarrow{b'}\widehat T'[\ell],
\]
where $b'$ is a homogeneous operation.  Consequently all
cokernel components of its overlap map vanish.
In degree $i\equiv r+1\pmod D$, the only kernel component is
the one belonging to this summand, and $\partial_r$ identifies
it with $\mathcal K$.  Thus $b'$ annihilates the projected
subspaces in \eqref{obs:killed-pair-bound}.
If the relevant operation degree is not a multiple of $D$,
then $b'=0$ already.  Otherwise normalize by the corresponding
power of $t$ and regard its coefficient, still denoted by $b'$,
as an element of $(\Delta[x^{-1}])^\Lambda$ acting on $K_0$.

We claim that $b'\in Sa$.  If not, apply the second part of
Lemma~\ref{obs:right-annihilator} to the span of all the
projected subspaces just found.  It produces an integral
operator $b_1$ annihilating them with
$\sigma(\bar b_1)\notin(C)$.  The vectors are also killed by
$a$, lie in $R$, and at scale $N$ contain a family independent
modulo $x^{qN+1}R$ of dimension
$cN^\beta-O(N^{1/3})$.  Lemma~\ref{obs:growth} contradicts
this because $\beta>1/3$.  The claim follows.

After extension to $\mathscr B$, a left multiple of $a$ kills
the cofiber connecting map, since $a[1]\partial_r=0$.
Hence every $\theta\widehat g$ is zero after this extension,
on every summand of $Y$.
For a perfect right $\mathscr E$-module $M$, evaluation gives a
natural equivalence of left $\mathscr B$-modules
\[
 \mathscr B\otimes_{\mathscr E}^{\mathbf L}M^*
 \xrightarrow{\ \sim\ }(M_{\mathscr B})^*.
\]
It is an equivalence for the free module and its shifts.  Both
sides are exact contravariant functors, so the class of modules
for which it is an equivalence is closed under finite sums,
cones and retracts.  Thus it holds for every perfect module.
Right flatness of $\mathscr B$ over $\mathscr E$ in
\eqref{obs:coefficient-extension} then gives
\[
 H\bigl((M_{\mathscr B})^*\bigr)
 \cong\mathscr B\otimes_{\mathscr E}H(M^*).
\]
For the module representing $\widehat P$, each homogeneous class
on the right is a finite sum of tensors $b_\nu\otimes[\theta_\nu]$.
Here each $[\theta_\nu]$ is represented, under strict realization,
by a pre-extension map to a shift of $\widehat T'$ as above.
By naturality, the dual of $\widehat g_{\mathscr B}$ sends this tensor
to $b_\nu\cdot[(\theta_\nu)_{\mathscr B}\widehat g_{\mathscr B}]$.
Every such composite has just been shown to vanish.  This proves
the right side of \eqref{obs:ghost-implication}.
\end{proof}

\subsection{Applying the criterion to the glued object}

Proposition~\ref{obs:dual-ghost} supplies the constraint on global
maps required by the finite-tower criterion.  The nonzero torsion
homology was established in \eqref{rew:obs:source-homology}, so
we can now conclude the obstruction.

\begin{theorem}
\label{obs:nonmembership}
For the object $U\in\mathbf D(\mathcal N)$ constructed in
Proposition~\ref{glue:object},
\[
 U\notin\Thick_{\mathbf D(\mathcal N)}(T').
\]
\end{theorem}

\begin{proof}
Take $S_0=S$, $D_0=D=2(p-1)$ and $\mathscr B_0=\mathscr B$ in
Theorem~\ref{rew:obs:tower-criterion}.  Lemma~\ref{obs:principal}
and the positive even value of $D$ verify its ring and period
assumptions.  Use the full thick subcategory
$\mathcal C_{\mathrm{comp}}$ and the exact functor $\Psi$ defined
above.  The subcategory contains $U$ and $T'$, and
$\Psi(T')\simeq\mathscr B$.  Since
$\mathcal C_{\mathrm{comp}}$ is full and thick in
$\mathbf D(\mathcal N)$ and contains $T'$, the ordinary thick
closure of $T'$ computed in either category is the same.
Equation~\eqref{rew:obs:source-homology} and
Lemma~\ref{obs:principal} verify that $H(\Psi(U))$ is torsion
in every degree and nonzero.  Proposition~\ref{obs:dual-ghost}
verifies \eqref{rew:obs:two-tests} for every map into that thick
closure, including every tower tail.  Theorem~\ref{rew:obs:tower-criterion}
therefore gives
\[
 U\notin\Thick_{\mathcal C_{\mathrm{comp}}}(T')
   =\Thick_{\mathbf D(\mathcal N)}(T').
\]
\end{proof}

Theorem~\ref{obs:nonmembership}, together with
Proposition~\ref{glue:object}, establishes all three conditions
of the transfer criterion: degreewise finite homotopy, a finite
complete construction, and failure of a global finite construction.
The remaining step is to transfer this object to the claimed
connective spectrum and record its chromatic consequences.
